\documentclass[a4paper,12pt]{article}
\usepackage[utf8]{inputenc}
\usepackage[T1]{fontenc}
\usepackage[ngerman]{babel}
\usepackage{amsmath, amssymb}
\usepackage{geometry}
\usepackage{tikz}
\usetikzlibrary{shapes.geometric, arrows.meta, positioning, fit, backgrounds}
\usepackage{parskip}
\usepackage{titlesec}
\usepackage{booktabs}
\usepackage{enumitem}

\geometry{left=3cm, right=3cm, top=2.5cm, bottom=2.5cm}

\titleformat{\section}{\large\bfseries}{}{0em}{}[\titlerule]
\titleformat{\subsection}{\normalsize\bfseries}{}{0em}{}

\tikzset{
  box/.style={rectangle, draw, rounded corners=3pt, text width=4.2cm,
              align=center, minimum height=0.9cm, font=\small, inner sep=4pt},
  boxw/.style={box, text width=5.2cm},
  boxn/.style={box, text width=3.8cm},
  arrow/.style={-{Latex[length=2mm]}, thick},
}

\begin{document}

\begin{center}
  {\LARGE\bfseries Studierhinweise zu Lektion 5}\\[0.5em]
  {\large Analysis (Modul 61211)}
\end{center}

\vspace{1em}

Diese Kurseinheit behandelt die Integralrechnung für reellwertige Funktionen auf
Intervallen. Der erste Abschnitt frischt die Kenntnisse aus dem Kurs MG über
Riemann-Unter- und Obersummen, Integrierbarkeit, Hauptsatz, partielle Integration
und Substitutionsregel auf; sorgfältiger sollten der Mittelwertsatz der
Integralrechnung und die Integrierbarkeit der Grenzfunktion einer gleichmäßig
konvergenten Funktionenfolge durchgearbeitet werden.

Der zweite Abschnitt erweitert den Integralbegriff auf \emph{uneigentliche
Integrale}, sodass auch unbeschränkten Funktionen oder Funktionen auf
unbeschränkten Intervallen ein Integral zugeordnet werden kann. Im dritten
Abschnitt werden \emph{Parameterintegrale} studiert und Bedingungen für
Stetigkeit und Differenzierbarkeit angegeben (Leibnizsche Regel). Der vierte
Abschnitt führt zu den \emph{Fourierreihen} $2\pi$-periodischer Funktionen, zur
Besselschen Ungleichung und zum Darstellungssatz. Der fünfte Abschnitt
kulminiert im \emph{Weierstraßschen Approximationssatz}, der zeigt, dass jede
auf einem kompakten Intervall stetige Funktion gleichmäßig durch
Polynomfunktionen approximierbar ist.

% -----------------------------------------------------------------------
\section*{Struktur der Kurseinheit 5}

\begin{center}
\begin{tikzpicture}[node distance=0.5cm and 0.6cm]

  % Column headers
  \node[font=\bfseries\small] (hA) at (0,0)    {Vorkenntnisse};
  \node[font=\bfseries\small] (hB) at (9.5,0)  {Neue Inhalte};

  % ---- Left column: Vorkenntnisse ----
  \node[box, below=0.8cm of hA] (v1) {Rückblick auf MG\\und Ergänzungen};
  \node[box, below=0.4cm of v1] (v2) {5.1 Riemannintegral,\\Rechenregeln};
  \node[box, below=0.4cm of v2] (v3) {5.1 Hauptsatz der\\Differenzial- und\\Integralrechnung};
  \node[box, below=0.4cm of v3] (v4) {5.1 Mittelwertsatz der\\Integralrechnung};
  \node[box, below=0.8cm of v4] (v5) {3.1 Grenzwerte von\\Funktionen,\\Rechenregeln};
  \node[box, below=0.4cm of v5] (v6) {2.3 Stetigkeits-,\\3.3/3.5 Differenzier-\\barkeitsbegriff};
  \node[box, below=0.4cm of v6] (v7) {2.6 Konvergenz von\\Funktionenreihen};

  % ---- Right column: new content ----
  \node[box, below=0.8cm of hB] (n1) {5.2 uneigentliches\\Integral};
  \node[box, below=0.4cm of n1] (n2) {5.2 Rechenregeln für\\uneigentliche Integrale};
  \node[box, below=0.4cm of n2] (n3) {5.2 Konvergenzkriterien\\(Cauchy, Majoranten)};
  \node[box, below=0.4cm of n3] (n4) {5.2 Beispiel:\\Gammafunktion};
  \node[box, below=0.4cm of n4] (n5) {5.3 Parameterintegral,\\Stetigkeit};
  \node[box, below=0.4cm of n5] (n6) {5.3 Leibnizsche Regel,\\Differenzierbarkeit};
  \node[box, below=0.4cm of n6] (n7) {5.4 $2\pi$-periodische\\Funktion, Fourierreihe};
  \node[box, below=0.4cm of n7] (n8) {5.4 Darstellbarkeit\\durch Fourierreihen};
  \node[box, below=0.4cm of n8] (n9) {5.5 Satz von Fejér};
  \node[box, below=0.4cm of n9] (n10){5.5 Approximationssatz\\von Weierstraß};

  % Arrows left column
  \draw[arrow] (v1) -- (v2);
  \draw[arrow] (v2) -- (v3);
  \draw[arrow] (v3) -- (v4);
  \draw[arrow] (v4) -- ++(0,-0.6) -- (v5.north);
  \draw[arrow] (v5) -- (v6);
  \draw[arrow] (v6) -- (v7);

  % Arrows right column
  \draw[arrow] (n1) -- (n2);
  \draw[arrow] (n2) -- (n3);
  \draw[arrow] (n3) -- (n4);
  \draw[arrow] (n4) -- (n5);
  \draw[arrow] (n5) -- (n6);
  \draw[arrow] (n6) -- (n7);
  \draw[arrow] (n7) -- (n8);
  \draw[arrow] (n8) -- (n9);
  \draw[arrow] (n9) -- (n10);

  % Cross arrows (left -> right)
  \draw[arrow, dashed] (v4.east) -- ++(0.3,0) |- (n1.west);
  \draw[arrow, dashed] (v5.east) -- ++(0.3,0) |- (n3.west);
  \draw[arrow, dashed] (v6.east) -- ++(0.3,0) |- (n6.west);
  \draw[arrow, dashed] (v7.east) -- ++(0.3,0) |- (n7.west);

\end{tikzpicture}
\end{center}

\newpage

% -----------------------------------------------------------------------
\section*{Zielelement 5.1 -- Rückblick und Ergänzungen: Das Riemannintegral}

\subsection*{Lerninhalte}

\begin{center}
\begin{tikzpicture}[node distance=0.5cm and 1.0cm]
  \node[box] (par)  {Partition von $I$,\\Unter-/Obersumme};
  \node[box, below=0.5cm of par]   (rie) {Riemann-integrierbare\\Funktion ($\mathcal{R}(I)$),\\Riemannintegral};
  \node[box, below=0.5cm of rie]   (eig) {Eigenschaften des\\Riemannintegrals\\(Additivität, Monotonie, \ldots)};
  \node[box, right=0.8cm of eig]   (mon) {monotone Funktion\\ist integrierbar};
  \node[box, below=0.5cm of eig]   (ubi) {unbestimmtes Integral,\\Stammfunktion};
  \node[box, below=0.5cm of ubi]   (hsz) {Hauptsatz der Differenzial-\\und Integralrechnung};
  \node[box, below=0.5cm of hsz]   (psi) {partielle Integration,\\Substitutionsregel};
  \node[box, below=0.5cm of psi]   (mws) {Mittelwertsatz der\\Integralrechnung};
  \node[box, below=0.5cm of mws]   (grf) {Integrierbarkeit der\\Grenzfunktion,\\gliedweise Integration};
  \draw[arrow] (par) -- (rie);
  \draw[arrow] (rie) -- (eig);
  \draw[arrow] (mon.south) |- (ubi.east);
  \draw[arrow] (eig) -- (ubi);
  \draw[arrow] (ubi) -- (hsz);
  \draw[arrow] (hsz) -- (psi);
  \draw[arrow] (psi) -- (mws);
  \draw[arrow] (mws) -- (grf);
\end{tikzpicture}
\end{center}

\subsection*{Lernziele}

Nach Durcharbeiten dieses Abschnitts sollten Sie sich aus MG wieder
vergegenwärtigt haben,
\begin{itemize}
  \item wie das Riemannsche Integral einer Funktion $f: [a,b] \to \mathbb{R}$ definiert
        ist und welche Regeln für das Rechnen mit ihm gelten,
  \item wie Differenzial- und Integralrechnung durch den ,,Hauptsatz``
        miteinander verknüpft sind,
  \item wie der Hauptsatz und die daraus abgeleiteten Regeln der partiellen
        Integration und der Integration durch Substitution anzuwenden sind.
\end{itemize}
Außerdem sollten Sie in der Lage sein,
\begin{itemize}
  \item den Mittelwertsatz der Integralrechnung zu formulieren und mit dem
        Mittelwertsatz der Differenzialrechnung in Verbindung zu bringen,
  \item den Satz über die Integrierbarkeit der Grenzfunktion einer gleichmäßig
        konvergenten Folge integrierbarer Funktionen zu formulieren.
\end{itemize}

\subsection*{Selbstkontrollelement 5.1}

Können Sie mittels $\dfrac{2}{1-x^2} = \dfrac{1}{1+x} + \dfrac{1}{1-x}$ das
Integral $\displaystyle\int_a^b \dfrac{1}{1-x^2}\,dx$ für $-1 < a < b < 1$ berechnen?

\newpage

% -----------------------------------------------------------------------
\section*{Zielelement 5.2 -- Uneigentliche Integrale}

\subsection*{Lerninhalte}

\begin{center}
\begin{tikzpicture}[node distance=0.5cm and 1.0cm]
  \node[box] (uni)  {uneigentliches Integral\\$\int_a^b f(x)\,dx$,\\Konvergenz};
  \node[box, below=0.5cm of uni]   (ver) {Verträglichkeit mit\\eigentlichem Integral};
  \node[box, below=0.5cm of ver]   (kri) {Cauchykriterium,\\Majorantenkriterium};
  \node[box, right=0.8cm of kri]   (rec) {Rechenregeln für\\uneigentliche Integrale};
  \node[box, below=0.5cm of kri]   (gam) {Beispiel: Gammafunktion\\$\Gamma(x) = \int_0^\infty t^{x-1}e^{-t}\,dt$};
  \node[box, below=0.5cm of gam]   (fun) {Funktionalgleichung\\$\Gamma(x+1) = x\,\Gamma(x)$,\\$\Gamma(n+1) = n!$};
  \draw[arrow] (uni) -- (ver);
  \draw[arrow] (ver) -- (kri);
  \draw[arrow] (rec.south) |- (gam.east);
  \draw[arrow] (kri) -- (gam);
  \draw[arrow] (gam) -- (fun);
\end{tikzpicture}
\end{center}

\subsection*{Lernziele}

Nach Durcharbeiten dieses Abschnitts sollten Sie in der Lage sein,
\begin{itemize}
  \item den Begriff des uneigentlichen Integrals zu definieren und an Beispielen
        zu erläutern,
  \item die Gammafunktion zu definieren und ihre Funktionalgleichung anzugeben.
\end{itemize}

\subsection*{Selbstkontrollelement 5.2}

,,Sehen`` Sie, warum das Integral $\displaystyle\int_1^\infty t^\alpha\,dt$ für
$\alpha < -1$ konvergiert, für $\alpha \geq -1$ dagegen nicht?

\newpage

% -----------------------------------------------------------------------
\section*{Zielelement 5.3 -- Parameterintegrale}

\subsection*{Lerninhalte}

\begin{center}
\begin{tikzpicture}[node distance=0.5cm and 1.0cm]
  \node[box] (par)  {Parameterintegral\\$F(x) = \int_\alpha^\beta K(x,t)\,dt$};
  \node[box, right=0.8cm of par]   (upi) {uneigentliches\\Parameterintegral};
  \node[box, below=0.5cm of par]   (stp) {Stetigkeit des\\Parameterintegrals};
  \node[box, right=0.8cm of stp]   (stu) {Stetigkeit des\\uneigentlichen\\Parameterintegrals};
  \node[box, below=0.5cm of stp]   (dif) {Differenzierbarkeit\\des Parameterintegrals\\(Leibnizsche Regel)};
  \node[box, right=0.8cm of dif]   (dui) {Differenzierbarkeit\\des uneigentlichen\\Parameterintegrals};
  \draw[arrow] (par) -- (upi);
  \draw[arrow] (par) -- (stp);
  \draw[arrow] (upi) -- (stu);
  \draw[arrow] (stp) -- (stu);
  \draw[arrow] (stp) -- (dif);
  \draw[arrow] (stu) -- (dui);
  \draw[arrow] (dif) -- (dui);
\end{tikzpicture}
\end{center}

\subsection*{Lernziele}

Nach Durcharbeiten dieses Abschnitts sollten Sie in der Lage sein, für
Funktionen, die als Parameterintegrale gegeben sind,
\begin{itemize}
  \item Bedingungen für ihre Stetigkeit und für ihre Differenzierbarkeit
        anzugeben,
  \item die Gammafunktion zu definieren und ihre Funktionalgleichung anzugeben.
\end{itemize}

\subsection*{Selbstkontrollelement 5.3}

Sei $F: \mathbb{R} \to \mathbb{R}$ durch $F(x) := \displaystyle\int_0^1 e^{-xt}\,dt$
gegeben. Berechnen Sie $F'(0)$
\begin{itemize}
  \item durch Auswerten des Integrals für $F(x)$,
  \item mithilfe der Leibnizschen Regel.
\end{itemize}

\newpage

% -----------------------------------------------------------------------
\section*{Zielelement 5.4 -- Fourierreihen}

\subsection*{Lerninhalte}

\begin{center}
\begin{tikzpicture}[node distance=0.5cm and 1.0cm]
  \node[box] (per)  {$2\pi$-periodische Funktionen,\\$C_k, S_k$,\\trigonometrische Reihe};
  \node[box, below=0.5cm of per]   (ska) {Skalarprodukt $\langle f,g\rangle$,\\Orthogonalitätsrelationen};
  \node[box, below=0.5cm of ska]   (fou) {Fourierreihe von $f$,\\Fourierkoeffizienten\\(Euler-Fouriersche Formeln)};
  \node[box, below=0.5cm of fou]   (bes) {Besselsche Ungleichung,\\$a_k, b_k \to 0$};
  \node[box, below=0.5cm of bes]   (dir) {Teilsummen $s_n$ und\\Dirichletkern $D_n$};
  \node[box, below=0.5cm of dir]   (dar) {Darstellungssatz:\\einseitige Differenzierbarkeit\\$\Rightarrow s_n(x_0) \to f(x_0)$};
  \draw[arrow] (per) -- (ska);
  \draw[arrow] (ska) -- (fou);
  \draw[arrow] (fou) -- (bes);
  \draw[arrow] (bes) -- (dir);
  \draw[arrow] (dir) -- (dar);
\end{tikzpicture}
\end{center}

\subsection*{Lernziele}

Nach Durcharbeiten dieses Abschnitts sollten Sie in der Lage sein,
\begin{itemize}
  \item die Fourierreihe einer $2\pi$-periodischen Funktion $f$, die über
        $]-\pi, \pi]$ integrierbar ist, aufzustellen,
  \item ein hinreichendes Kriterium dafür anzugeben (und in Beispielen
        anzuwenden), dass die Funktion in einem Punkt durch ihre Fourierreihe
        dargestellt wird.
\end{itemize}

\subsection*{Selbstkontrollelement 5.4}

Es sollte für Sie kein Problem sein zu begründen, dass die Fourierreihe jeder
$2\pi$-periodischen, differenzierbaren Funktion $f: \mathbb{R} \to \mathbb{R}$
gebildet werden kann und dass die Fourierreihe für jedes $x \in \mathbb{R}$
gegen $f(x)$ konvergiert.

\newpage

% -----------------------------------------------------------------------
\section*{Zielelement 5.5 -- Der Weierstraßsche Approximationssatz}

\subsection*{Lerninhalte}

\begin{center}
\begin{tikzpicture}[node distance=0.5cm and 1.0cm]
  \node[box] (fej)  {Fejérsche Summen $\sigma_n$,\\Fejérkerne $F_n$};
  \node[box, below=0.5cm of fej]   (eig) {Eigenschaften der $F_n$:\\nichtnegativ, gerade,\\$\frac{1}{\pi}\int_{-\pi}^\pi F_n\,du = 1$};
  \node[box, below=0.5cm of eig]   (sfe) {Satz von Fejér:\\gleichmäßige Approximation\\$2\pi$-period.\ stetiger Fkt.\\durch $\sigma_n$};
  \node[box, below=0.5cm of sfe]   (apw) {Approximationssatz\\von Weierstraß:\\stetige $f: [a,b] \to \mathbb{R}$\\durch Polynomfunktionen};
  \draw[arrow] (fej) -- (eig);
  \draw[arrow] (eig) -- (sfe);
  \draw[arrow] (sfe) -- (apw);
\end{tikzpicture}
\end{center}

\subsection*{Lernziele}

Nach Durcharbeiten dieses Abschnitts sollten Sie in der Lage sein,
\begin{itemize}
  \item den Weierstraßschen Approximationssatz zu formulieren.
\end{itemize}

\subsection*{Selbstkontrollelement 5.5}

Seien $a, b \in \mathbb{R}$, $a < b$. Der Beweis der folgenden Aussage sollte
Ihnen keine Mühe machen:
\[
  f \in C([a,b]) \; \Rightarrow \; \forall\, \varepsilon > 0\; \exists\, P \in \mathcal{P}\;
  \forall\, x \in [a,b]: |f(x) - P(x)| < \varepsilon.
\]

\end{document}
